\documentclass[11pt]{article}
\usepackage[utf8]{inputenc}
\usepackage{amsmath,amssymb}
\usepackage[margin=1in]{geometry}
\usepackage{hyperref}
\emergencystretch=3em

\title{Complete expansions of the chart integrals with arbitrary $(k,h)$:\\
$Z(\sqrt n)=n^{-p/2}KD\sum_{i\le j}c_{j,i}\big(\tfrac12\log n+S\log b\big)^i+O(n^{-p/2-\varepsilon/2})$}
\author{Claude, with Daniel Murfet}
\date{2026-09-06}

\begin{document}
\maketitle
\sloppy

\begin{abstract}
Units~50 and~52 reduced the chart integrals exactly to the form $Z(\sqrt n)=K(\sqrt n)^{-p}D\,
H_j(\sqrt n\,b^S)$ for a divide-and-integrate tower $H$ on an admissible base with tail rate
$\varepsilon$. This unit turns the base-independent expansion of units~49/51 into a single generic
theorem: $Z(\sqrt n)=n^{-p/2}KD\sum_{i\le j}c_{j,i}(\tfrac12\log n+S\log b)^i+O(n^{-p/2-\varepsilon/2})$,
instantiated for the all-equal case with arbitrary $(k,h)$ ($\varepsilon=1$) and the mixed case
($\varepsilon=q-p$). A bridge lemma identifies unit~39's \texttt{iterChart} with \texttt{iterChartGen}
at $k=1$, $h=0$. Zero \texttt{sorry}/\texttt{axiom}.
\end{abstract}

\section{Setting}
With $L=\sqrt n\,b^S$ one has $\log L=\tfrac12\log n+S\log b$ and
$L^{-\varepsilon}=b^{-S\varepsilon}n^{-\varepsilon/2}$; the expansion
$|H_j(L)-\sum c_{j,i}\log^iL|\le(M/\varepsilon^j)L^{-\varepsilon}$ therefore transfers to an explicit
$O(n^{-p/2-\varepsilon/2})$ remainder once multiplied by the prefactor $K(\sqrt n)^{-p}D$.

\section{The things formalised}
{\small
\begin{verbatim}
noncomputable def towerCoeff (G : ℝ → ℝ) (A : ℝ) (j : ℕ) (b : ℝ) (S : ℕ) (n : ℝ) : ℝ :=
  ∑ i ∈ Finset.range (j + 1),
    divCoeff G A j i * (Real.log n / 2 + (S : ℝ) * Real.log b) ^ i
theorem tower_expansion_isBigO (hG : LogBase G A M C δ ε) (j : ℕ) (K D p b : ℝ) (S : ℕ)
    (hb : 0 < b) (Z : ℝ → ℝ)
    (hZ : ∀ n, 0 < n →
      Z n = K * (Real.sqrt n) ^ (-p) * D * iterDivPrim G j (Real.sqrt n * b ^ S)) :
    (fun n => Z n - n ^ (-(p / 2)) * (K * D * towerCoeff G A j b S n))
      =O[atTop] fun n => n ^ (-(p / 2) - ε / 2)
theorem iterChartGen_expansion_isBigO ...        -- ε = 1, remainder n^(-(p/2) - 1/2)
theorem iterChartGen_mixed_expansion_isBigO ...  -- ε = q - p
theorem iterChart_eq_iterChartGen (β a b : ℝ) (d : ℕ) (c : ℝ) :
    iterChart β a b d c = iterChartGen β a b (fun _ => 1) (fun _ => 0) d c
\end{verbatim}
}

\section{Proof strategy}
The generic theorem is \texttt{IsBigO.of\_bound} with the constant $|KD|\,(M/\varepsilon^j)\,
b^{-S\varepsilon}$, valid for $n\ge\max(1,b^{-2S})$ so that $L\ge1$; the three \texttt{rpow}
identities ($(\sqrt n)^{-p}=n^{-p/2}$, $(\sqrt n)^{-\varepsilon}=n^{-\varepsilon/2}$,
$(b^S)^{-\varepsilon}=b^{-S\varepsilon}$) are isolated as \texttt{have}s. The two instances supply
only the exact-reduction hypothesis \texttt{hZ} from units~50 and~52. The bridge lemma is an
induction with \texttt{pow\_zero}, \texttt{pow\_one}.

\section{Roadblocks and resolutions}
\begin{itemize}
\item A \texttt{ring} after a \texttt{rw} that had already closed the goal (\texttt{log\_pow}
produced the exact form); remove it. Otherwise first-pass.
\end{itemize}

\section{What was Mathlib, what was new}
Mathlib: \texttt{Real.rpow\_add}, \texttt{Real.rpow\_mul}, \texttt{Real.sqrt\_eq\_rpow},
\texttt{Real.log\_pow}. New: the generic tower expansion and its two instances.

\section{Lessons}
Once the exact reductions are stated in a common shape, the asymptotic and expansion theorems become
one generic lemma plus one-line instances; stating the reductions with that shape in mind (constants
$K,D$, scale $b^S$) is worth the small awkwardness at the reduction stage.

\section{Outcome}
For arbitrary $(k,h)$ the chart integrals of \texttt{thm:TaylorTree} in the all-equal and one-noncritical
regimes now have complete polynomial-in-$\log n$ expansions with explicit recursive coefficients and
power-saving remainders. Remaining at this level: several noncritical coordinates, and the product-set
(Fubini) form of the $d$-fold integral.
\end{document}
